\subsection{IT Governance}

\indent 

\noindent \textbf{IT Governance}: the framework of \textbf{leadership, structures, and processes} that ensures an organisation’s \textbf{IT systems and strategies support its overall business goals}, deliver value, and manage risks.

\subsubsection{Why IT Governance matters in Professional Practice}

\indent 

\noindent\textbf{Accountability and Responsibility}: 

IT professionals don’t just “build systems” → they must ensure systems are secure, reliable, and aligned with business needs. Governance frameworks clarify who is accountable (e.g., CIO, project leads, developers).

\noindent\textbf{Ethical and Professional standards}:

Poor governance can lead to privacy breaches, unethical data use, or financial loss. As professionals, you’re bound by codes of conduct (ACS, IEEE, ISACA) to act responsibly and escalate risks.

\noindent\textbf{Risk awareness}

IT projects often fail due to poor governance (scope creep, lack of oversight). Professionals must recognise governance as part of risk management and compliance (e.g., GDPR, APRA CPS 234).

\noindent\textbf{Alignment with Business Strategy}

In practice, IT is not about “cool tech” alone → it’s about delivering measurable business value. Governance ensures your work supports organisational goals, not just technical outputs.

\noindent\textbf{Career relevance}

Many IT roles (data analyst, project manager, architect, engineer) require understanding of governance frameworks like COBIT, ITIL, ISO 38500. Demonstrating governance awareness makes you more valuable as a professional.

\subsubsection{Understanding IT Governance Success}

\indent 

IT Governance is considered \textbf{successful} when an organisation’s IT \textbf{delivers value, manages risk, and aligns with business strategy}.

\noindent\textbf{Key elements of IT Governance success}:

\begin{itemize}
    \item \textbf{Strategic alignment}: IT projects and initiatives directly support the organisation’s goals
    \item \textbf{Value delivery}: IT investments generate measurable business benefits (cost savings, revenue, improved services).
    \item \textbf{Risk management}: Cybersecurity, data privacy, operational risks are identified, monitored, and reduced.
    \item \textbf{Resource optimisation}: IT resources (people, infrastructure, applications) are used effectively and not wasted.
    \item \textbf{Performance measurement}: KPIs, Balanced Scorecards, and regular audits ensure IT performance is tracked and improved
\end{itemize}

\noindent\textbf{Indicators of Success}

\begin{itemize}
    \item IT projects delivered on time, on budget, and with expected business outcomes.
    \item High levels of executive support and stakeholder satisfaction.
    \item Compliance with laws and regulations.
    \item Evidence that IT decisions are transparent, ethical, and accountable
\end{itemize}

\subsubsection{Why Governance Fails?}

\noindent\textbf{Lack of executive support}

Governance requires buy-in from top management (CIO, CEO, Board). Without leadership backing, policies are ignored, and accountability is weak.

\noindent\textbf{Misalignment with business goals}

IT initiatives operate in isolation from business strategy. Result: Projects may be technically sound but deliver little value.

\noindent\textbf{Poor communication and awareness}

Employees and managers don’t understand governance rules or processes. Leads to inconsistent implementation and confusion.

\noindent\textbf{Resistance to change}
Governance often requires process changes, documentation, or oversight. Teams may resist if seen as bureaucratic or slowing delivery.

\subsubsection{Governance Frameworks}

\noindent\textbf{IT Governance Framework}: a \textbf{structured set of practices}, \textbf{principles}, and \textbf{processes} that guide how IT is managed to support business strategy, deliver value, and control risks.

\noindent\textbf{Why do we need Frameworks?}

\begin{itemize}
    \item Provide a common language between IT and business.
    \item Define roles, responsibilities, and accountability.
    \item Ensure compliance with laws and regulations.
    \item Improve efficiency and performance measurement.
\end{itemize}

\subsection{COBIT}

\indent 

\noindent \textbf{C}ontrol \textbf{O}bjectives for Information and related \textbf{T}echnologies.

\noindent\textbf{COBIT}: A \textbf{globally recognised IT governance and management framework} developed by \textbf{ISACA}. Provides \textbf{principles, practices, models, and tools} for IT governance.

\bigskip

It is structured into 2 parts:

\begin{itemize}
    \item \textbf{Governance Objectives}: Evaluate, Direct, Monitor (EDM).
    \item \textbf{Management Objectives}: Plan, Build, Run, Monitor (PBRM).
\end{itemize}

\subsubsection{Objectives of COBIT}

\noindent\textbf{Governance Objectives}

Governance objectives are about \textbf{directing and overseeing IT}. It asks "Are we doing the right things? These objectives are usually owned by \textbf{boards and executives} (CIO, steering committees, board of directors).

They follow the EDM model:

\begin{itemize}
    \item \textbf{Evaluate}: Assess stakeholder needs, business goals, risks, and compliance requirements.
    \item \textbf{Direct}: Set priorities, make governance decisions, allocate resources.
    \item \textbf{Monitor}: Continuously check performance, compliance, and stakeholder satisfaction
\end{itemize}

\bigskip

\noindent\textbf{Management Objectives}

Management objectives are about the \textbf{day-to-day running of IT}. It asks "Are we doing thing the right way?" These objectives are usually owned by \textbf{IT managers, project managers, and operational teams}.

They follow the PBRM model:

\begin{itemize}
    \item \textbf{Plan}: Strategy, budget, enterprise architecture, project planning.
    \item \textbf{Build}: Develop solutions, acquire infrastructure, configure systems.
    \item \textbf{Run}: Operate services, service desk, incident management, change management.
    \item \textbf{Run}: Monitor: Evaluate performance, security, service availability, risk monitoring.
\end{itemize}

\subsubsection{Evolution of COBIT}

\indent 

\step COBIT 1 - Focused on control objectives for IT auditing.

\step COBIT 2 - Broader IT management and governance elements added.

\step COBIT 3 - Alignment with business goals, risk management focus.

\step COBIT 4 - Comprehensive governance + management, integrated other frameworks (e.g., ITIL, ISO).

\step COBIT 5 – (Five Key Governance Principles)

\step COBIT 2019 (6th Principle)

\subsubsection{COBIT 5: Five Key Governance Principles}

\begin{itemize}
    \item \textbf{Meeting Stakeholder Needs}: Ensures IT goals align with business objectives and stakeholder expectations.
    \item \textbf{Covering the Enterprise End-to-End}: Governance is applied across the entire organization, not just the IT department.
    \item \textbf{Applying a Single Integrated Framework}: COBIT 5 integrates with other standards (e.g., ITIL, ISO/IEC 27001, TOGAF).
    \item \textbf{Enabling a Holistic Approach}: Governance is driven by interconnected enablers: processes, principles, culture, information, infrastructure, and people.
    \item \textbf{Separating Governance from Management}: Governance sets direction and evaluates, while management plans, builds, and executes.
\end{itemize}

\subsubsection{COBIT 2019: The Sixth Principle}

\begin{itemize}
    \item \textbf{Tailoring Governance to Enterprise Needs}: COBIT 2019 introduced \textbf{design factors} that allow governance to be customized based on the enterprise’s unique strategy, risk profile, size, regulatory environment, and priorities.
\end{itemize}

\subsubsection{What changed from COBIT 5 to COBIT 2019?}

\begin{itemize}
    \item \textbf{Expanded governance principles}: COBIT 5 featured 5 key governance principles; COBIT 2019 adds a sixth, updating the framework's foundation.
    \item \textbf{Increased process count}: The framework expanded from 37 processes in COBIT 5 to 40 in COBIT 2019, introducing updates such as “Managed Assurance,” “Managed Projects,” and “Managed Data.” 
    \item \textbf{Shift from “Enablers” to “Components”}: COBIT 5’s “enablers” were renamed “components” in COBIT 2019, giving a broader and more flexible structural backdrop for designing governance systems. 
    \item \textbf{Open-source-style model}: COBIT 2019 incorporates feedback from practitioners and is periodically updated under the oversight of a steering committee, ensuring continuous improvement and relevance.
\end{itemize}

\subsubsection{Purpose of COBIT}

\begin{itemize}
    \item \textbf{Strategic Alignment}: Ensure IT supports and drives business objectives.
    \item \textbf{Value Delivery}: Ensure IT delivers benefits against the strategy, concentrating on optimising costs and resources.
    \item \textbf{Risk Management}: Identify, manage, and mitigate IT risks.
    \item \textbf{Resource Optimisation}: Maximise efficient use of people, processes, and technology.
    \item \textbf{Performance Measurement}: Provide metrics and maturity models to measure IT effectiveness.
\end{itemize}

\subsection{ITSM - ITIL v4.0}

\indent

\noindent\textbf{IT Service Management (ITSM)} is the \textbf{discipline and practice of managing IT as a set of services} that deliver value to customers (internal or external).

\bigskip

ITSM focuses on:

\begin{itemize}
    \item Processes (how IT is delivered)
    \item People (who uses and supports IT)
    \item Services (the outcomes IT provides to the business)
\end{itemize}

\subsubsection{COBIT vs ITIL (ITSM)}

\indent 

\noindent\textbf{COBIT}: A \textbf{governance and management framework} for enterprise IT. Focuses on \textbf{what should be done} to align IT with business goals. Owned by \textbf{ISACA}.

\noindent\textbf{ITIL}: A \textbf{best practice framework for IT service management (ITSM)}. Focuses on \textbf{how IT services should be delivered and managed}. Owned by \textbf{AXELOS} (a joint venture of UK Cabinet Office \& Capita).

\subsubsection{History of ITSM}

\noindent 1960s

IT was mainly hardware-focused (mainframes, punch cards, batch processing). IT departments were technology-centric, not service-oriented. Support was ad-hoc, with no formal processes. Users often had to “know someone in IT” to get help.

\bigskip

\noindent 1980s
Organisations realised IT needed to be reliable and standardised. Government agencies (especially in the UK) wanted consistent IT services across departments. The UK government’s CCTA (Central Computer and Telecommunications Agency) created ITIL v1 (late 1980s) $\xrightarrow{}$ a collection of best practices for managing IT.

\bigskip

\noindent 1990s

ITIL v2 (1990s) simplified processes into Service Support and Service Delivery. It became the global standard for ITSM best practices. Widely adopted by industries like banking, telecom, and government. ISO/IEC 20000 (2005) emerged as the first international standard for ITSM, based heavily on ITIL.

\bigskip 

\noindent 2000s

ITIL v3 (2007) introduced the Service Lifecycle model. ITSM grew beyond IT support into business-aligned IT delivery. Frameworks like COBIT, MOF (Microsoft Operations Framework), and ISO standards reinforced ITSM practices.

\bigskip

\noindent 2010s

ITIL 2011 refined ITIL v3, making processes clearer. ITIL 4 (2019) integrated Agile principles → making ITSM more flexible and adaptive to digital transformation. ITSM shifted from rigid processes to value co-creation with customers.

\bigskip

\noindent Today

ITSM now focuses on Cloud services (AWS, Azure, SaaS), Automation \& AI(chatbots, self-service portals, predictive analytics), and User experience (UX)as a core outcome. ITSM is no longer just “IT operations”, it is about enabling digital business strategy.

\subsubsection{ITIL v3 Process}

\noindent\textbf{Service Strategy}

The main purpose of this step is to define\textbf{ IT as a strategic service provider} that delivers value. The main focus is on Business outcomes, value creation, service portfolio, financial management, risk management.

\noindent\textbf{Service Design}

The main purpose of this step is to design IT services and processes to meet business needs. The main focus is on Service catalog, availability, capacity, security, continuity, and compliance.

\noindent\textbf{Service Transition}

The main purpose of this step is to \textbf{Build, test, and deploy IT services} safely. The main focus is on Change management, release management, knowledge management, service validation.

\noindent\textbf{Service Operation}

The main purpose of this step is to Run and support IT services efficiently. The main focus is on Incident management, problem management, request fulfillment, monitoring, service desk.

\noindent\textbf{Continual Service Improvement}

The main purpose of this step is to continuously measure and improve services based on feedback and performance metrics. The main focus is on KPI monitoring, service reviews, lessons learned, process improvements.

\subsubsection{How ITIL Processes work together}

\begin{itemize}
    \item Incident occurs $\xrightarrow{}$ Incident Management
    \item Root cause identified $\xrightarrow{}$ Problem Management
    \item Change required $\xrightarrow{}$ Change Enablement
    \item Update deployed $\xrightarrow{}$ Deployment Management
    \item Service measured $\xrightarrow{}$ Continual Improvement
\end{itemize}

ITIL processes are \textbf{end-to-end service lifecycle activities}, ensuring IT services are delivered efficiently, reliably, and aligned with business needs.

\subsubsection{ITSM, COBIT and ITIL}

\begin{table}[h]
\centering
\begin{tabular}{|c|l|l|l|}
\hline
\textbf{Feature} & \multicolumn{1}{c|}{\textbf{ITSM}} & \multicolumn{1}{c|}{\textbf{COBIT}} & \multicolumn{1}{c|}{\textbf{ITIL}} \\ \hline
\textbf{Type}    & Discipline / Practice              & Governance Framework                & Best Practice Framework            \\ \hline
\textbf{Focus}   & Delivering IT as a service         & Alignment, control, value, risk     & Implementing ITSM processes        \\ \hline
\textbf{Scope}   & Service management                 & Governance + management             & Service lifecycle                  \\ \hline
\textbf{Level}   & Operational \& Tactical            & Strategic \& Tactical               & Operational \& Tactical            \\ \hline
\end{tabular}
\end{table}




\subsection{ITIL v4}

\noindent

\textbf{ITIL v4} is \textbf{the latest version of the IT Service Management (ITSM) framework}, released in 2019 by AXELOS. It is \textbf{value-driven, flexible, and designed for modern digital environments}, including Agile, Lean, and DevOps.

Information Technology Infrastructure Library is the most widely used framework for IT Service Management (ITSM)

It shifts the focus from rigid processes (ITIL v3), to a \textbf{flexible, value-driven, and digital-first approach}.

ITIL is a set of practices. It's primary purpose is to provide a systematic approach to IT Service Management.

\subsubsection{Service Value System}

\indent 

The SVS ensures all components and activities work together to \textbf{co-create value}. It consists of 5 major components:

\begin{itemize}
    \item \textbf{Guiding Principles}: Universal recommendations for decision-making.
    \item \textbf{Governance}: Oversight and compliance.
    \item \textbf{Service Value Chain (SVC)}: Operating model for creating value.
    \item \textbf{Practices}: 34 ITIL practices replacing the older “processes.”
    \item \textbf{Continual Improvement}: Ongoing improvement of services and practices
\end{itemize}

\subsubsection{Guiding Principles}

\indent 

There are mainly 7 guiding principles:

\begin{itemize}
    \item \textbf{Focus on value}: All activities must deliver value to stakeholders.
    \item \textbf{Start where you are}: Use existing resources and processes effectively.
    \item \textbf{Progress iteratively with feedback}: Small, incremental improvements are better than large, risky changes.
    \item \textbf{Collaborate and promote visibility}: Ensure transparency and teamwork.
    \item \textbf{Think and work holistically}: Consider all aspects of the organization and service.
    \item \textbf{Keep it simple and practical}: Avoid unnecessary complexity.
    \item \textbf{Optimize and automate}: Streamline processes and leverage technology.
\end{itemize}

\subsubsection{Service Value Chain}

\indent 

This is the main operating model. The \textbf{operating model} shows how value is created through \textbf{end-to-end activities}.

There are mainly 6 key activities:

\begin{itemize}
    \item \textbf{Plan}: Ensure understanding of vision, current state, and direction.
    \item \textbf{Improve}: Continual improvement of services, practices, and processes.
    \item \textbf{Engage}: Interact with stakeholders to understand needs.
    \item \textbf{Design \& Transition}: Build and deploy new or changed services.
    \item \textbf{Obtain/Build}: Acquire or build service components.
    \item \textbf{Deliver \& Support}: Run services and support users.
\end{itemize}

\subsubsection{Practices}

\indent 

This replaces the ITIL v3 processes. There are a total of 34 practices noted in ITIL v4. These are broadly categorised into 3 major practices:

\begin{itemize}
    \item \textbf{General Management Practices (14)}: These include Strategy, Risk Management, Continual Improvement, Knowledge Management, Portfolio Management, etc.
    \item \textbf{Service Management Practices (17)}: These include Incident Management, Problem Management, Change Enablement, Service Desk, Service Level Management, Availability Management, etc.
    \item \textbf{Technical Management Practices (3)}: These include Deployment Management, Software Development \& Management, Infrastructure \& Platform Management
\end{itemize}

\subsection{Question}

\subsubsection{End of Lecture Questions}

\begin{itemize}
    \item Explain why IT Governance is critical for IT professionals.
    \item List three benefits of ITSM.
    \item Describe the difference between COBIT and ITIL.
    \item Name the seven guiding principles of ITIL v4.
    \item Briefly explain the Service Value Chain activities in ITIL v4.
\end{itemize}

\subsubsection{Case Study}

\indent 

The University of Sydney wants to implement a new Learning Management System (LMS) to improve student engagement and faculty collaboration. The university’s IT team is responsible for delivering the LMS, ensuring it aligns with strategic goals, is reliable, and meets compliance standards.

Additional Context:

Past IT projects at the university experienced delays and budget overruns.

Some staff resist adopting new digital tools due to unfamiliarity.

The IT team plans to use COBIT 2019 for governance, ITIL v4 for ITSM, and agile practices for development.

\bigskip

\noindent\textbf{Question}:

\begin{itemize}
    \item Using COBIT, what governance practices should the IT team implement to ensure project alignment with university goals?
    \item Which ITIL v4 Service Value Chain activities would the IT team apply to ensure smooth LMS deployment?
    \item Suggest three ITIL v4 practices that could improve adoption among faculty and students.
    \item How does integrating COBIT governance with ITIL v4 practices help co-create value in this scenario?
\end{itemize}


